\documentclass[12pt,a4paper,openany,oneside]{book}

% --- Paquetes ---
\usepackage[utf8]{inputenc}
\usepackage[spanish, es-noshorthands]{babel}
\usepackage{amsmath, amsfonts, amssymb}
\usepackage{graphicx}
\usepackage{geometry}
\usepackage{hyperref}
\usepackage{booktabs}
\usepackage{fancyhdr}
\usepackage{listings}
\usepackage{xcolor}
\usepackage{tikz}
\usepackage{pgfplots}
\usepackage{colortbl}
\usepackage{multirow}
\usepackage{hhline}
\usetikzlibrary{arrows.meta, positioning, shapes.geometric, calc, shapes, chains, scopes, shadows}
\pgfplotsset{compat=1.18}

\geometry{margin=2.5cm}

% --- Colores y estilos para código Python ---
\definecolor{codegreen}{rgb}{0,0.6,0}
\definecolor{codegray}{rgb}{0.5,0.5,0.5}
\definecolor{codepurple}{rgb}{0.58,0,0.82}
\definecolor{backcolour}{rgb}{0.95,0.95,0.92}

\lstset{
    language=Python,
    backgroundcolor=\color{backcolour},   
    commentstyle=\color{codegreen},
    keywordstyle=\color{blue},
    numberstyle=\tiny\color{codegray},
    stringstyle=\color{codepurple},
    basicstyle=\ttfamily\small,
    breaklines=true,
    frame=single,
    showstringspaces=false
}

% --- Estilo de cabecera y pie ---
\pagestyle{fancy}
\fancyhf{}
\renewcommand{\headrulewidth}{0.4pt}
\renewcommand{\footrulewidth}{0.4pt}
\fancyhead[L]{\small Carlos César Sánchez Coronel}
\fancyhead[R]{\small Sesión 12: RAG y Bases de Datos Vectoriales}
\fancyfoot[L]{\small Carlos César Sánchez Coronel}
\fancyfoot[C]{\thepage}

% --- Portada ---
\title{
    \textbf{Sesión 12} \\ 
    \Huge \textbf{RETRIEVAL-AUGMENTED GENERATION (RAG)} \\
    \Large Y Bases de Datos Vectoriales
}
\author{
    \small Por \\ \vspace{0.2cm}
    \Large \textsc{Carlos César Sánchez Coronel}
}
\date{2026}

\begin{document}

\maketitle
\tableofcontents
\addcontentsline{toc}{chapter}{Índice general}

% ------------------------------------------------------------------
% CAPÍTULO 12: RAG Y BASES DE DATOS VECTORIALES
% ------------------------------------------------------------------
\chapter{Retrieval-Augmented Generation (RAG) y Bases de Datos Vectoriales}

En la era de los Large Language Models (LLMs), uno de los mayores desafíos es el manejo de información privada o dinámica que no formó parte del entrenamiento original del modelo. Esta sesión explora la arquitectura RAG, una técnica que permite a los modelos "consultar" una base de conocimientos externa antes de generar una respuesta, reduciendo drásticamente las alucinaciones y permitiendo la actualización de conocimientos sin necesidad de reentrenamiento costoso.

\section{Logro de la sesión}
Al finalizar la sesión, el estudiante será capaz de diseñar e implementar sistemas que combinan la recuperación de información (Retrieval) con la generación de lenguaje (Generation), utilizando bases de datos vectoriales para realizar búsquedas semánticas eficientes.

\section{Limitaciones de los LLMs}

A pesar de su capacidad, los LLMs sufren de dos problemas críticos que RAG busca resolver:

\begin{enumerate}
    \item \textbf{Conocimiento Desactualizado (Knowledge Cutoff)}: Los modelos están limitados a los datos con los que fueron entrenados. Si un modelo terminó su entrenamiento en 2024, no sabrá eventos de 2025.
    \item \textbf{Alucinaciones}: Cuando un modelo no tiene la información, tiende a generar respuestas que suenan coherentes pero son fácticamente falsas.
    \item \textbf{Falta de Datos Privados}: Un modelo público (como GPT-4) no conoce los documentos internos de una empresa específica.
\end{enumerate}

\section{¿Qué es RAG (Retrieval-Augmented Generation)?}

RAG es un framework que optimiza la salida de un LLM al consultar una base de conocimientos confiable fuera de sus datos de entrenamiento iniciales antes de generar una respuesta.

\subsection{El Flujo de Trabajo RAG}
El proceso se divide en tres etapas principales:
\begin{itemize}
    \item \textbf{Recuperación (Retrieval)}: El sistema busca documentos relevantes en una base de datos basándose en la consulta del usuario.
    \item \textbf{Aumento (Augmentation)}: La consulta original y los documentos recuperados se combinan en un "Prompt" enriquecido.
    \item \textbf{Generación (Generation)}: El LLM utiliza este contexto para redactar la respuesta final.
\end{itemize}

\begin{center}
\begin{tikzpicture}[
    node distance=1.5cm,
    block/.style={rectangle, draw, fill=blue!10, text width=2.5cm, textAlign=center, minimum height=1.2cm, rounded corners, font=\small, align=center, drop shadow},
    arrow/.style={-Stealth, thick}
]
    % Nodos
    \node[block] (user) {Usuario (Pregunta)};
    \node[block, right=of user] (search) {Búsqueda Semántica};
    \node[block, below=of search, fill=green!10] (db) {Base de Datos Vectorial};
    \node[block, right=of search] (augment) {Prompt Enriquecido (Pregunta + Contexto)};
    \node[block, right=of augment, fill=orange!10] (llm) {LLM (Generador)};
    \node[block, below=of llm] (resp) {Respuesta Final};

    % Flechas
    \draw[arrow] (user) -- (search);
    \draw[arrow] (search) -- (augment);
    \draw[arrow] (db) -- (search);
    \draw[arrow] (augment) -- (llm);
    \draw[arrow] (llm) -- (resp);

\end{tikzpicture}
\captionof{figure}{Diagrama del flujo de un sistema RAG estándar.}
\end{center}

\section{Embeddings y Búsqueda Semántica}

Para que el sistema "entienda" qué documentos son relevantes, no buscamos palabras clave (como en Google antiguo), sino **significado**. 

\subsection{Embeddings}
Son representaciones numéricas (vectores de alta dimensión) de fragmentos de texto. Textos con significados similares terminan en puntos cercanos dentro del espacio vectorial.

\[ \text{Vector}(\text{"Rey"}) - \text{Vector}(\text{"Hombre"}) + \text{Vector}(\text{"Mujer"}) \approx \text{Vector}(\text{"Reina"}) \]

\subsection{Similitud del Coseno}
Es la métrica más usada para medir qué tan cerca están dos vectores $\mathbf{A}$ y $\mathbf{B}$:
\[ \text{similarity} = \cos(\theta) = \frac{\mathbf{A} \cdot \mathbf{B}}{\|\mathbf{A}\| \|\mathbf{B}\|} \]

\section{Bases de Datos Vectoriales}

A diferencia de las bases de datos SQL tradicionales, las vectoriales están optimizadas para almacenar y buscar entre millones de vectores de forma casi instantánea.

\begin{table}[h]
\centering
\begin{tabular}{@{}lll@{}}
\toprule
\textbf{Herramienta} & \textbf{Tipo} & \textbf{Uso Principal} \\ \midrule
\textbf{FAISS} & Librería (Meta) & Búsqueda local, alta velocidad en RAM. \\
\textbf{Pinecone} & Managed Service & Escalabilidad en la nube, nivel empresarial. \\
\textbf{Weaviate} & Open Source & Base de datos vectorial con búsqueda híbrida. \\
\textbf{ChromaDB} & Open Source & Fácil integración con cuadernos de Python. \\ \bottomrule
\end{tabular}
\caption{Principales tecnologías de almacenamiento vectorial.}
\end{table}

\section{Evaluación de Sistemas RAG: El "RAG Triad"}

Evaluar RAG es difícil porque hay dos partes que pueden fallar. Usamos principalmente dos métricas de fidelidad:

\begin{itemize}
    \item \textbf{Faithfulness (Fidelidad)}: ¿La respuesta se deriva estrictamente del contexto recuperado o el modelo inventó cosas (alucinó)?
    \item \textbf{Relevance (Relevancia)}: ¿El contexto recuperado es realmente útil para responder la pregunta del usuario?
\end{itemize}

\section{Implementación Conceptual (Python)}

\begin{lstlisting}[language=Python]
# Ejemplo simplificado usando LangChain y OpenAI
from langchain_community.vectorstores import FAISS
from langchain_openai import OpenAIEmbeddings, ChatOpenAI
from langchain.chains import RetrievalQA

# 1. Preparar datos y embeddings
embeddings = OpenAIEmbeddings()
vector_db = FAISS.from_texts(["La Sesion 12 trata sobre RAG", "Carlos es el profesor"], embeddings)

# 2. Configurar el LLM
llm = ChatOpenAI(model_name="gpt-4", temperature=0)

# 3. Crear la cadena RAG
rag_chain = RetrievalQA.from_chain_type(
    llm=llm, 
    chain_type="stuff", 
    retriever=vector_db.as_retriever()
)

# 4. Consulta
query = "De que trata la Sesion 12?"
response = rag_chain.invoke(query)
print(response["result"])
\end{lstlisting}

\section{Preparación para Entrevistas y Aplicaciones}

\subsection{Pregunta Clave: "¿Qué es RAG y cuándo lo usarías?"}
\textbf{Respuesta}: RAG es una arquitectura que combina la recuperación de documentos con modelos generativos. Lo usaría cuando:
\begin{itemize}
    \item Se requiere trabajar con datos privados que no pueden subirse a la nube para entrenamiento.
    \item La información cambia frecuentemente (noticias, precios, inventarios).
    \item Es crítico reducir las alucinaciones citando fuentes específicas.
\end{itemize}

\subsection{Diferencia entre RAG y Fine-Tuning}
\begin{center}
\begin{tabular}{|l|p{5cm}|p{5cm}|}
\hline
\rowcolor{gray!20} \textbf{Característica} & \textbf{Fine-Tuning} & \textbf{RAG} \\ \hline
\textbf{Conocimiento} & Estático (en los pesos). & Dinámico (en la base de datos). \\ \hline
\textbf{Costo} & Alto (requiere GPUs y tiempo). & Bajo (solo costo de inferencia/DB). \\ \hline
\textbf{Alucinaciones} & Posibles. & Muy reducidas. \\ \hline
\textbf{Transparencia} & Caja negra. & Puede citar fuentes (links/docs). \\ \hline
\end{tabular}
\end{center}

\section{Resumen}
\begin{itemize}
    \item \textbf{RAG} soluciona el problema de conocimiento desactualizado y alucinaciones.
    \item Los \textbf{Embeddings} convierten texto en vectores para búsqueda por significado.
    \item Las \textbf{Bases de Datos Vectoriales} (FAISS, Pinecone) son el motor que almacena el conocimiento externo.
    \item \textbf{Faithfulness} y \textbf{Relevance} son claves para asegurar que el sistema no mienta.
\end{itemize}

\vspace{1cm}
\begin{center}
    \textit{"RAG no es solo una técnica, es el puente entre la inteligencia del modelo y la realidad de los datos actuales."}
\end{center}

\end{document}